\chapter{Prostatitis Symptoms}

\section*{Definition}
Prostatitis refers to inflammation of the prostate gland, presenting with varying combinations of pelvic pain, urinary symptoms, and sexual dysfunction. It is classified into four NIH categories: I (acute bacterial), II (chronic bacterial), III (chronic prostatitis/chronic pelvic pain syndrome -- CP/CPPS), and IV (asymptomatic inflammatory).

\section*{What Happens in Your Body}
Bacterial prostatitis involves ascending infection from the urethra, reflux of infected urine into prostatic ducts, or hematogenous spread. CP/CPPS (Category III, >90\% of cases) involves a complex interplay of pelvic floor hypertonicity, neurogenic inflammation, central sensitization, and autoimmune/immune signaling proteins-mediated inflammation without active infection.

\section*{Causes}
\begin{itemize}
  \item \textbf{Acute bacterial (I):} E. coli (most common, ~80\%), Klebsiella, Proteus, Pseudomonas, Enterococcus
  \item \textbf{Chronic bacterial (II):} Same organisms, often with biofilm formation in prostatic calculi
  \item \textbf{Chronic pelvic pain syndrome (III):} Multifactorial -- pelvic floor muscle spasm, repetitive trauma (cycling), psychological stress, depression, anxiety, autoimmune trigger, interstitial cystitis, pudendal neuralgia
  \item \textbf{Risk factors:} Indwelling catheter, urinary tract instrumentation, unprotected receptive anal intercourse, anatomic abnormalities, BPH, recent UTI, immunosuppression
  \item \textbf{Sexually transmitted:} Neisseria gonorrhoeae, Chlamydia trachomatis (uncommon but possible)
\end{itemize}

\section*{When to See a Doctor}
See your doctor if:
\begin{itemize}
  \item Fever, chills, and severe pelvic pain with urinary retention (suspect acute bacterial prostatitis -- may be febrile)
  \item Chronic pelvic pain persisting >3 months with urinary or ejaculatory symptoms
  \item Hematuria, hematospermia, or systemic symptoms (weight loss, bone pain)
\end{itemize}

\section*{Related Symptoms}
\begin{itemize}
  \item Perineal pain, suprapubic pain, dysuria, urinary frequency, urgency, hesitancy, weak stream, post-void dribbling, ejaculatory pain, hematospermia, erectile dysfunction, low libido
\end{itemize}
\textbf{Important:} CP/CPPS is a diagnosis of exclusion; a thorough evaluation for other causes of pelvic pain is essential.

\section*{Self-Care / Home Management}
\begin{itemize}
  \item Sit in a warm sitz bath for 15--20 minutes to relax pelvic muscles
  \item Avoid prolonged sitting, cycling, and heavy lifting
  \item Increase fluid intake to encourage frequent urination
  \item Reduce or eliminate caffeine, alcohol, and spicy foods
\end{itemize}

\section*{How Your Doctor Will Evaluate You}
\begin{itemize}
  \item History and physical including digital rectal exam (DRE): tender, swollen, warm prostate in acute prostatitis
  \item Urinalysis and culture with pre- and post-massage urine (Meares--Stamey 4-glass test) for categories I/II
  \item Serum PSA (may be elevated in prostatitis; do not measure in acute phase)
  \item In chronic (III): NIH Chronic Prostatitis Symptom Index (NIH-CPSI)
  \item Ultrasound or MRI to exclude abscess, stones, or malignancy in atypical cases
\end{itemize}

\section*{Treatment Options}
\begin{itemize}
  \item \textbf{Acute bacterial (I):} Empiric broad-spectrum antibiotics (e.g., fluoroquinolone or third-gen cephalosporin) then tailored by culture; hospitalization if systemically ill or urinary retention
  \item \textbf{Chronic bacterial (II):} Long-course (4--12 weeks) fluoroquinolone or trimethoprim-sulfamethoxazole with good prostatic penetration
  \item \textbf{CP/CPPS (III):} Multimodal approach: alpha-blockers (tamsulosin), anti-inflammatories (NSAIDs), pelvic floor physical therapy, and CBT
  \item \textbf{For all types:} Alpha-blockers for urinary symptoms, stool softeners, NSAIDs for pain
\end{itemize}

\section*{References}
\begin{thebibliography}{99}
\bibitem{prostatitis_symptoms:ref1} Krieger JN, Nyberg L Jr, Nickel JC. "NIH consensus definition and classification of prostatitis." \textit{JAMA}, 1999;282(3):236--237.
\bibitem{prostatitis_symptoms:ref2} Nickel JC, Shoskes DA, Wagenlehner FME. "Management of chronic prostatitis/chronic pelvic pain syndrome." \textit{European Urology}, 2013;63(1):128--140.
\bibitem{prostatitis_symptoms:ref3} Schaeffer AJ. "Clinical practice: Chronic prostatitis and the chronic pelvic pain syndrome." \textit{New England Journal of Medicine}, 2006;355(16):1690--1698.
\bibitem{prostatitis_symptoms:ref4} Rees J, Abrahams M, Doble A, et al. "Diagnosis and treatment of chronic bacterial prostatitis in adults." \textit{BMJ}, 2015;351:h4796.
\bibitem{prostatitis_symptoms:ref5} Anothaisintawee T, Attia J, Nickel JC, et al. "Management of chronic prostatitis/chronic pelvic pain syndrome: A systematic review." \textit{JAMA}, 2011;305(1):78--86.
\bibitem{prostatitis_symptoms:ref6} Wagenlehner FME, Naber KG, Bschleipfer T, et al. "Prostatitis and male pelvic pain syndrome: Diagnosis and treatment." \textit{Deutsches \"{A}rzteblatt International}, 2018;115(4):55--60.
\end{thebibliography}
